\documentclass[11pt,a4paper]{article}
\usepackage[T1]{fontenc}
\usepackage{amsmath,amssymb,geometry}
\geometry{margin=27mm}
\title{EPPA--III Section 3: exceptional \(C_5\) and \(R_3\) validation}
\author{Independent Lean validation supplement}
\date{16 October 2026}
\begin{document}
\maketitle
\noindent\textbf{Policy.} The author manuscript
\texttt{main-new(7).tex} is frozen. This note records
proposed local proof replacements and formalisation
boundaries; no manuscript text is edited.

\section*{Lemma 3.18: neighbourhood-pair orbit argument}
Only the cross-degree-two case is addressed.
Write \(G=C_5\) and \(B=V(H)\setminus V(G)\), so
\(|B|=5\), and assume \(|N_H(w)\cap V(G)|=2\) for
every \(w\in B\). The cyclic order-five subgroup of
\(\operatorname{Aut}(G)\) has two orbits on two-element
subsets of \(V(G)\), each of size five:
the edges and the diagonals of \(C_5\).
Every base automorphism lifts to an automorphism of
\(H\) preserving \(B\), so the multiplicity
\(m(S)=|\{w\in B:N_H(w)\cap V(G)=S\}|\) is constant
along these orbits.
As \(\sum_{|S|=2}m(S)=5\), the two orbit multiplicities
are \(1,0\) in one order or the other.
Thus the profiles are pairwise distinct and precisely
one orbit occurs. This excludes a \(C_4\dot\cup C_6\)
cross-incidence, which necessarily repeats a profile.
Unlike the original prime-part-size shortcut, the
argument uses the \emph{actual EPPA profile multiset},
not a presumed action on an indexed outside half.

\paragraph{Lean coverage.}
The modules \texttt{CycleFivePairProfileOrbits.lean},
\texttt{CycleFiveOutsideProfileSaturation.lean}, and
\texttt{CycleFiveOutsideEquiv.lean} establish, under
these hypotheses, the two pair orbits, their multiplicity
saturation, a genuine equivalence
\(V(C_5)\simeq B\), and injectivity of the outside profile
map. \texttt{CycleFiveInvariantFourForms.lean}
independently proves that a five-vertex graph invariant
under all cyclic shifts is empty, \(C_5\), its
complement, or complete.

\paragraph{Still outstanding.}
The classification of \(H[B]\) under its
\emph{actual} induced action, edge-count elimination
and any remaining cross-degree branches of Lemma~3.18
are not claimed by these modules. The complete
Lemma~3.18 is not yet certified.

\section*{Lemma 3.20: canonical degree-four rook conclusion}
Let \(R_3\) denote the \(3\times3\) rook graph.
Suppose \(e:R_3\hookrightarrow H\) is an
induced EPPA embedding, \(|H|=18\), \(H\) is
vertex-transitive, and
\[
  |N_H(w)\cap e(V(R_3))|=4
  \qquad(w\in H\setminus e(V(R_3))).
\]
The orbit-multiplicity argument leaves exactly one
of the two nine-element orbits of 4-subsets, the
rook open neighbourhoods or strict nonneighbourhoods.
Each profile is realised once, giving an \emph{actual}
equivalence from rook vertices to \(B\).
Genuine EPPA lifts show that the complementary
graph is invariant under the rook automorphisms,
hence has one of four forms. Vertex-transitivity
and equal internal edge counts eliminate empty
and complete forms. The complementary-rook
top graph contradicts equality of triangle counts
across the two halves; the two remaining graphs
are the permitted doubles.

\paragraph{Exact Lean theorem.}
\texttt{rook\_eppa\_degree\_four\_host\_is\_rook\_double}
produces an equivalence
\(\phi:\{0,1\}\times V(R_3)\simeq V(H)\)
fixing the designated bottom copy and identifying
the whole graph with one of
\(\texttt{rookDoubleGraph false crossNonrook}\).
It does not assume a pre-indexed second half,
coherent lifts or a triangle-count formula.

\paragraph{Still outstanding.}
The theorem assumes the outside cross-degree is
exactly four. To complete Lemma~3.20 one must
derive or classify the other possible degrees
under the \emph{actual} hypotheses, then connect
the explicit rook doubles to the author's named
alternatives. The homogeneous-host and final
Proposition~3.23 cases are not implied by this
conditional theorem.

\paragraph{Continuation: actual \(C_5\) outside adjacency.}
The new \texttt{CycleFiveOutsideAction.lean} theorem
\texttt{cycleFive\_eppa\_degree\_two\_outside\_two\_forms}
combines the saturated orbit labels with genuine
EPPA transport of rotations, then uses
vertex-transitivity and the equality of internal
edge counts to show that the \emph{actual}
complementary induced graph is \(C_5\) or its
complement. The identification is through an
outside equivalence \emph{constructed} from the
profiles, not assumed at the outset.
Pending consolidated Lean CI and axiom audit,
this is a proposed certification milestone.
To finish the cross-degree-two case, exclude
the complementary-top configuration when it
is incompatible with host transitivity (the
triangle-count route), and identify the
surviving actual host graphs. It is not a
proof of all cross-degree cases.

\section*{Minimal source-facing edits}
For Lemma~3.18, replace only the claim that
the prime size of the bipartition forces one
cycle by the two-orbit multiset count above.
For Lemma~3.20, replace the assumed
simultaneous action on both halves by the
profile-multiplicity argument before identifying
the canonical doubles. See the shorter existing
\texttt{validation/short-r3-cross.tex}.
Retain the author's prose elsewhere pending review.

\paragraph{Certification rule.}
Count the named Lean declarations as verified only
after the consolidated GitHub Actions build and
\texttt{CheckAxioms.lean} audit pass. The earlier
independent successful PR runs do not automatically
validate imports in combination.
\end{document}
